\section{Specialized accelerators and learning}\label{sec:accelerators}

\subsection{A native propositional-search lane}

Djinn's logical-search heritage remains useful. Implement a native
proof-producing propositional solver as a specialized lane, using a focused or
contraction-controlled calculus for its precisely reified fragment. The
reification records how every connective and atom maps to an exact Lean type
or proposition. A positive derivation is translated into a Lean term and
checked. This can be much cheaper than unconstrained enumeration for pure
higher-order plumbing.

Do not reuse a propositional decision procedure as a universal decision
procedure for dependent inhabitation. A case where the reification deliberately
abstracts a dependent family as an atom may still produce a valid positive
term after reconstruction, but failure of the abstraction is not a full-Lean
emptiness result. The fragment certificate must say what was represented and
what was omitted. This preserves the distinction between the predecessor's
logical lane and its broader bounded synthesis extensions~\cite{djex-readme}.

Some polymorphic impossibilities have direct Lean certificates. For example,
an alleged value of
\code{forall A B : Type, A -> B} can be specialized to an inhabited source and
an empty target to derive a contradiction. That is a genuine semantic
refutation in the stated context. It should be distinguished from both bounded
search failure and non-derivability of an unrelated propositional abstraction.

\subsection{Finite synthesis and arithmetic templates}

Small finite domains admit an additional exact route: enumerate or solve a
finite truth table, reify it as a Lean expression, and prove coverage of the
domain. For bounded bitvector adapters and finite state transitions, a SAT/SMT
choice grammar can be effective. The relevant completeness bound includes the
permitted expression grammar and domain enumeration, not just the solver's
Boolean search.

For arithmetic functions, use typed templates such as affine expressions,
piecewise affine expressions with a bounded number of guards, and polynomial
forms of bounded degree. Solver results instantiate coefficients; Lean proves
the resulting universal specification in the supported theory. A found formula
can be useful even when the grammar is incomplete. A failed coefficient search
cannot exclude a non-template solution.

\subsection{Optional learned proposal ranking}

A learned model can rank providers, suggest generalized motives, or propose
small sketches. Its output should enter the same typed-plan interface as any
other heuristic rule. It cannot change the request's policy, add a hypothesis,
or turn a raw prediction into a proof.

Training examples should include failed partial branches and their validated
failure reasons, not just successful final programs. The controller can learn
which obligation order reduces branching or which invariant grammar to try
first. Keep a reproducible nonlearned reference lane and report ablations;
otherwise a change in the model silently changes the effective synthesis
language. No model is required for the initial Leant2 implementation.

\section{Evaluation: measuring practical coverage honestly}\label{sec:evaluation}

\subsection{Benchmark requests, not screenshots}

Each benchmark records the original type, local context, exact universes,
contract, imports, allowed term and proof providers, semantic policy, and
resource bounds. Preserve binder order and dictionary selections. A easier
variant that succeeds is a useful diagnostic but remains a different benchmark.

Use at least four complementary strata: provider-free structural tasks;
realistic library-adapter tasks; dependent/indexed construction; and
behavior-sensitive recursive programs. Include proofs and law-carrying
structures without allowing the relevant completed target declaration itself
to become a provider. Keep effectful and noncomputable profiles in separate
strata because their success criteria differ.

Hold out projects or families of related types, not merely random rows from a
corpus containing near-duplicate signatures. A solution that memorizes the
standard term for a familiar Church signature should not be mistaken for
coverage of arbitrary dependent data. Conversely, legitimate library reuse is
valuable and should be measured in its own stratum rather than prohibited in
all evaluations.

\subsection{Metrics and controls}

Report the fraction of original requests with verified implementations, not
the fraction for which some type-correct term was emitted. Separate finite
example contracts from universal contracts; separate supplied sketches from
fully synthesized programs. Record first verified result, explored alternatives,
proof time, search time, peak memory, generated-code complexity, and the selected
assumption profile. Report medians and tails, not only the best examples.

The most informative controls exercise boundaries:
\begin{itemize}
\item Wrong total programs at the exact type, with actual observed/proved
rejection; false contracts with evidence that the predicate check ran.
\item Alpha-renaming, binder permutation with appropriately transformed types,
fresh universe names, and two distinguishable dictionaries of one class type.
\item Metavariable assignment leakage, stale environment generations, invalid
cache reuse after a contract change, and source re-elaboration under changed
instances.
\item Impossible indexed branches, nondefinitional arithmetic casts, and
attempted recursion without a decrease capability.
\item Solver unknown, malformed certificates, worker crashes, cancelled
requests, and late replies from obsolete solver frames.
\end{itemize}

Ablations should disable one mechanism at a time: continuation-aware
backtracking, result-head indexing, contract propagation, proof-carrying
motives, length abstractions, retained observation classes, or learned ranking.
Keep the same original request and budget. The finite experiment already
provides two such controls, but production conclusions require the wider
benchmark suite.

\subsection{Metamorphic and generated tests}

Generate well-typed terms in a restricted grammar, erase selected subterms to
make synthesis requests, and check whether the declared bound includes the
known witness. This tests relative coverage without requiring a human oracle
for every case. Transform the requests by alpha-renaming, universe renaming,
wrapping in equivalent structures, or adding irrelevant locals. Changes in
performance are permissible; changes in verified meaning are not.

Use deliberately different-but-same-typed values to test evidence identity.
For dependent terms, generate a first component and then the type of the second
component, rather than independently generating ill-typed pairs and blaming
the synthesizer for rejecting them. Keep generated counterexample input codecs
and proof of validity in the test artifact.

\section{Implementation sequence and acceptance gates}\label{sec:roadmap}

\begin{longtable}{@{}L{0.12\linewidth}L{0.34\linewidth}L{0.43\linewidth}@{}}
\caption{A Lean-only delivery sequence. Gates describe evidence, not dates.}\\
\toprule
Stage & Deliverable & Acceptance gate \\
\midrule
\endfirsthead
\toprule Stage & Deliverable & Acceptance gate \\\midrule
\endhead
\bottomrule\endfoot
M0 & Exact request/accepted-result protocol; LeanBridge; transaction tests; source export. & Fresh declaration checks reject unresolved holes, extra premises, stale handles, and forbidden dependencies. \\
M1 & Local and provider spine search; nested polymorphism; coupled goals; basic proof services; explicit scheduler. & Original-type replay for rank-N, independent universes, strict implicits, dictionaries, and contract-driven backtracking. \\
M2 & Generic constructors, dependent cases, transport search, indexed library retrieval. & User-defined indexed families and context-dependent elimination; no handwritten case skeleton counted as automatic discovery. \\
M3 & Recursor schemas, proof-carrying motives, contract propagation, typed CEGIS and first certified domains. & Provider-free map/append/reverse/filter-style tasks with universal contracts, retained alternatives, and valid negative controls. \\
M4 & Generalized motives, accumulator/measure synthesis, well-founded and selected mutual recursion; effect-specific adapters. & Original recursive requests pass program, contract, and termination checks; cost/assumption profiles remain explicit. \\
M5 & Wider project evaluation, optimized indexes/caches, optional SMT and learned ranking, robust editor/REPL integration. & Holdout benchmarks, failure accounting, bounded-memory operation, reproducible source export, and documented version support. \\
\end{longtable}

M0 and M1 should not depend on Z3, a learned model, or a Haskell service. The
microprototype is an experiment toward M1, not a delivery of M0's full
acceptance infrastructure. M2 and M3 provide the largest semantic gains over a
type-projection architecture: generic dependent construction and recursion
whose correctness obligations influence the program as it is built.

Porting the predecessor's whole module structure is not a prerequisite. Reuse
its test intent, acceptance distinctions, and useful logical algorithms; keep
its Haskell implementation as a historical comparison executable only when
needed for measurement. There is no architectural requirement to preserve a
Haskell frontend or a neutral Haskell-compatible public API in Leant2.

\subsection{The highest-risk implementation choices}

The first risk is \emph{state coupling}: an engine can produce individually
plausible subproofs while mishandling shared assignments. Address it with the
continuation invariant and adversarial tests before optimizing the queue.
The second is \emph{recursive generalization}: most useful recursive synthesis
failures will not be fixed by deeper blind enumeration. Invest in typed
motive, accumulator, and measure grammars with proof obligations.

The third is \emph{proof bottlenecks}. A correct program whose contract is
beyond the configured proof methods remains unresolved. Keep its candidate
and residual proof state, and distinguish that from program-construction
failure. The fourth is \emph{scaling exactness}: environment-aware caches,
provider retrieval, and branch snapshots must be measured under large imports
without weakening their identity contracts.

These are concrete engineering and research tasks, not reasons to weaken the
goal to a Haskell-shaped fragment. The architecture isolates them so that each
can be improved and evaluated without rewriting the semantic foundation.

\section{Conclusion}

Leant2 should synthesize \emph{typed programs with attached obligations}, not
strings that resemble inhabitants. Exact Lean expressions eliminate an
unnecessary semantic projection, while explicit search plans preserve control,
reproducibility, and scheduling. Dependent telescopes supply the right local
constraints; contracts distinguish inhabitants; recursors supply legal
recursion; theorem-backed abstractions eliminate impossible choices; and the
kernel checks the final program/proof pair.

The decisive implementation principle is that all these components constrain
the same object. A program hole, its dictionary choice, its index, its recursive
invariant, its behavioral observations, and its proof must not drift into
parallel approximations that can be recombined by name or printed text. The
continuation-aware Lean experiment and the observation-deduplication ablation
show two small but concrete consequences of this principle.

The recommended starting point is therefore a native exact-expression core
with strong transaction and acceptance boundaries, followed quickly by generic
dependent elimination and proof-carrying recursion. That gives Leant2 a route
to substantially broader practical Lean coverage without making universal
inhabitation, solver omniscience, or unmeasured performance promises part of its
contract.
